\section{Exact certificates for adaptive contraction plans}\label{sec:certificates}

Fix an exact mode-weight allocation. A plan partitions the $r$ group labels
into three ordered blocks. Inside each block, recursively split after its
floor midpoint and form the two boundary tables. Finish by contracting the
three complete block tables. A merge has parts $A,B,O$, where $O$ is the
complement of $A\cup B$. Its three index classes are the unordered edges
between those parts. The final operation has the three final blocks as
its parts. There are $r-2$ operations in a complete plan.

For the backend of Theorem~\ref{thm:updated-rectangle}, set
\[
 H_0=21\,400\,000.
\]
The nine integer support triples are the distinct permutations of
\begin{equation}\label{eq:supports}
 (21\,400\,107,21\,400\,107,0),\qquad
 (21\,400\,000,16\,750\,214,10\,000\,000).
\end{equation}
A support vector is zero on edges inside an operation part, and takes
the corresponding coefficient on each of its three cross classes.
Let $\mathcal V_j$ contain all $9(r-2)$ supports of plan $j$.
The local operation exponent of that plan is bounded by
\[
 F_j(g)=\frac1{H_0}\max_{v\in\mathcal V_j}v\cdot g.
\]
This includes every internal merge. It does not count only the final
triangle. The input/output branch in \eqref{eq:supports} also pays for
every stored boundary table.

\begin{lemma}[Weighted coupled-pair certificate]\label{lem:paircertificate}
Let there be $q\ge2$ plans, with fixed nonnegative integer weights
$w_1,\ldots,w_q$ and $W=\sum_jw_j>0$. For $i<j$ compute
\[
 A_{ij}=\max_{v\in\mathcal V_i,\,w\in\mathcal V_j}
 \left(w_i^2\norm v^2+w_j^2\norm w^2
             +2(q-1)w_iw_j(v\cdot w)\right).
\]
Choosing the least-cost plan gives local time and space coefficient
\begin{equation}\label{eq:betacertificate}
 \beta_r^2=
 \frac{2\sum_{i<j}A_{ij}}{(q-1)W^2H_0^2r^2}
\end{equation}
relative to $\K(M,u)$. For one plan the valid squared coefficient is
$2\max_{v\in\mathcal V_1}\norm v^2/(H_0^2r^2)$.
\end{lemma}

\begin{proof}
For a given $g\ge0$, choose $v_j$ attaining the maximum for plan $j$.
The exact identity
\[
 (q-1)\left\|\sum_jw_jv_j\right\|_2^2
 =\sum_{i<j}\left(w_i^2\norm{v_i}^2+w_j^2\norm{v_j}^2
                   +2(q-1)w_iw_jv_i\cdot v_j\right)
\]
holds because each diagonal term is counted $q-1$ times and each cross
term once with its required multiplier. Therefore the squared norm is
at most $\sum A_{ij}/(q-1)$. Also
\[
 \min_jF_j(g)\le\frac1{WH_0}\left(\sum_jw_jv_j\right)\cdot g.
\]
Apply \eqref{eq:unordered}. Only a fixed finite number of supports and
support sums is possible, so the inherited error is uniform. One need
not suppose the independently maximizing support pairs are jointly
consistent. An upper bound on their sum is sufficient. The single-plan
case follows directly from the same entropy estimate.

To choose a plan, write $d_{\{a,b\}}=|\mathcal I_{ab}|\,|\mathcal I_{ba}|$.
Compute for each plan the integer
\[
 P_j=\max_{v\in\mathcal V_j}\prod_e d_e^{v_e},
\]
and choose the smallest, breaking ties by its index. Its logarithm is
$H_0F_j(g)$. The number of supports, their coefficients and denominators
are fixed, so these integers have polynomial bit length. Selection
introduces no comparison oracle for real logarithms. All input/output
tables are covered by the first support family. Sequential matrix
execution then bounds local space by local time on the exponential scale.
\end{proof}

The identity is a finite mathematical tool. The new work consists of the
updated backend, the new fifteen-group construction, their exact
certificates and the paid global schedules. Averaging and quadratic
norm identities themselves are classical.

\subsection{The new fifteen-group cyclic construction}

Start with the ordered blocks
\begin{equation}\label{eq:15seed}
 (13,8,12,10,1)\mid(2,14,0,5,6)\mid(3,4,7,9,11).
\end{equation}
For $j=0,\ldots,14$, add $j$ modulo 15 to every entry, preserving the
within-block order. Use all fifteen resulting plans with equal weights.
Each five-element block is built by the prescribed recursive split, so
the internal order matters. There are 117 supports per plan.

Exact enumeration of all
$\binom{15}{2}117^2=1\,437\,345$ support pairs gives
\begin{equation}\label{eq:beta15}
 \boxed{\beta_{15}^2=
   \frac{10\,797\,642\,389\,573\,327}{38\,640\,375\,000\,000\,000}}.
\end{equation}
The generator, all plans, pair maxima and maximizing indices are supplied.
The checker uses arbitrary-precision integers throughout. An independent
implementation explicitly reconstructs the four merges inside each
five-element block and confirms the same fraction. This is an exhaustive
finite certificate, not a sample of entropy vectors or a numerical
optimization output. The seed was found by finite exploratory search;
no global optimality of the seed or schedule is claimed.

\subsection{Affine nine and cyclic thirteen}

The four parallel classes of the affine plane $\mathbb F_3^2$ give the
following useful complete plans:
\[
\begin{split}
 &(0,1,2)\mid(3,4,5)\mid(6,7,8),\\
 &(0,3,6)\mid(1,4,7)\mid(2,5,8),\\
 &(0,5,7)\mid(1,3,8)\mid(2,4,6),\\
 &(0,4,8)\mid(2,3,7)\mid(1,5,6).
\end{split}
\]
These plans are inherited from the latest affine packets. The updated
backend yields a new coefficient. Put
$a=335/428+10^{-5}$ and $b=50/107$. The exact value is
\begin{equation}\label{eq:beta9}
 \beta_9^2=
 \frac{3(1+a^2+b^2)+2(1+a+b)^2}{54}
 =\frac{357\,623\,761\,181\,449}{1\,236\,492\,000\,000\,000}.
\end{equation}
There is a geometric explanation for the expression: a final second-branch
support has squared norm $9H_0^2(1+a^2+b^2)$, and supports from transverse
directions have inner product $2H_0^2(1+a+b)^2$. Every intersection of
their three cross classes has exactly two edges. The largest coupled pair
value is $42\,914\,851\,341\,773\,880$. The exact checker exhausts
all $6\cdot63^2=23\,814$ pairs and checks that the internal-merge cases
do not exceed that value. All four final supports attain it jointly;
this certifies sharpness of this particular average-norm calculation,
not optimality for actual SAT instances or all plans.

For thirteen groups use all cyclic translates of the existing ordered seed
\[
 (11,2,8,9,7)\mid(10,1,0,5)\mid(6,4,3,12).
\]
With the new backend, the same exact method gives
\begin{equation}\label{eq:beta13}
 \beta_{13}^2=
 \frac{36\,196\,226\,703\,680\,947}{125\,767\,265\,000\,000\,000}.
\end{equation}
There are 99 supports per plan and $78\cdot99^2=764\,478$ checked pairs.

\subsection{Ten-group cache and small tails}

The ten-group stage retains the nine plans and mixing weights of
\cite{pairwise}. The weights, in plan order, are
\[
 (137,124,135,120,108,119,106,87,64),\qquad W=1000.
\]
The complete ordered plans are in \texttt{data/new\_backend\_weighted10.json}.
They are included rather than read from a mutable external location.
Applying Lemma~\ref{lem:paircertificate} to the new supports gives
\begin{equation}\label{eq:beta10}
 \beta_{10}^2=
 \frac{13\,753\,456\,054\,476\,639\,694\,573}
      {45\,796\,000\,000\,000\,000\,000\,000}.
\end{equation}
All $36\cdot72^2=186\,624$ pairs are checked. This does not require
the optional unary-potential reparameterization explored during the review.

The small-tail plans are $(0,1,2)\mid(3,4)\mid(5,6)$,
$(0,1)\mid(2,3)\mid(4,5)$, and $(0,1)\mid(2)\mid(3)$.
Their coefficients are
\[
 \beta_7^2=\frac{120001200003}{245000000000},\quad
 \beta_6^2=\frac{4444488889}{10000000000},\quad
 \beta_4^2=\frac{40000400001}{80000000000}.
\]
The input/output branch determines these bounds. Omitting it would
incorrectly shrink a cost that includes materializing the tables.

